\Legend
小明有 $n$ 件衣服，第 $i$ 件衣服的保暖值为 $a_i$。他要从中带走若干件，以应对接下来 $m$ 天的天气。每天可以从带走的衣服中任选一个\textbf{非空}子集穿着；不同天可以选择不同的子集。

第 $i$ 天给出 $t_i,x_i$：

\begin{itemize}
\item 如果 $t_i=1$，当天所穿衣服的保暖值之和必须\textbf{至少}为 $x_i$。
\item 如果 $t_i=2$，当天所穿衣服的保暖值之和必须\textbf{不超过} $x_i$。
\end{itemize}
求能满足所有天要求的最少携带件数。若不可能，输出 $-1$。

\Input

第一行整数 $n$，表示衣服数。第二行 $n$ 个整数 $a_1,\ldots,a_n$。第三行整数 $m$，表示天数。接下来 $m$ 行，每行两个整数 $t_i,x_i$。

保证 $1\le n\le3$，$1\le a_i\le2\times10^5$，$1\le m\le100$，$t_i\in\{1,2\}$，$1\le x_i\le10^9$。

\Output

输出一个整数，表示最少携带件数；无解时输出 $-1$。

\Examples

\begin{example}
\exmp{3
2 3 5
1
1 10}{3}
\end{example}

\begin{example}
\exmp{3
2 3 5
2
1 5
2 2}{2}
\end{example}

\begin{example}
\exmp{3
7 11 13
2
1 24
2 7}{3}
\end{example}

\begin{example}
\exmp{2
17 19
2
1 16
2 18}{1}
\end{example}

\begin{example}
\exmp{3
2 11 19
2
1 1000000000
2 1}{-1}
\end{example}
